\section{State-Access Ordering and Execution Consistency}
\label{app:revision_effects}

The revised executor makes the order of conflicting state accesses part of the
stored program. This supplies the missing ordering information when two nodes
have no data-dependency path but access the same external state. The original
data edges continue to carry method arguments; additional control dependencies
only determine when a node may execute.

\paragraph{Execution versions and evidence.}
We distinguish the original executor, denoted V0, from the revision with optional
effect ordering, denoted V1. The three-level type experiment, matched state
encoding experiment, CPU scaling measurement, and frozen-controller replay below
were completed with V0 and frozen protocols. Their measurements remain associated
with that implementation. V1 is evaluated separately by execution-consistency
tests and frozen-model regression. The subsequent pool and continual-learning
studies use their stated search procedures around the ATF representation and
execution interfaces. This separation identifies which implementation produced
each result while retaining the original experimental records.

\paragraph{Program order and conflict relation.}
For each method instance $v$, let $R_v$ and $W_v$ be its declared read and write
resource sets. Two instances conflict when
\begin{equation}
 W_u\cap(R_v\cup W_v)\ne\varnothing
 \quad\text{or}\quad
 W_v\cap(R_u\cup W_u)\ne\varnothing .
\label{eq:revision_conflict}
\end{equation}
Before the first evaluation of a candidate, a topological order $\tau$ of its
data graph is selected and stored using stable node identities. For each
conflicting pair, control dependencies preserve its order in $\tau$.
The implementation maintains the latest writer and outstanding readers for
each resource; transitive dependencies avoid storing every conflicting pair.
Unknown effects form conservative barriers among effectful methods, and a
wildcard resource denotes access to all resources. Read-only accesses and
accesses to disjoint resources retain their available parallelism.

All data and control dependencies point forward in $\tau$, so their union is
acyclic. Assuming deterministic methods, complete resource declarations, and
identical initial state and inputs, successful executions respecting this union
produce the same outputs and final state. To see this, any two legal linear
extensions are connected by adjacent swaps of incomparable nodes. Such nodes
have neither a data dependency nor a conflicting state access, so each swap
preserves their outputs and state updates. Repeating these swaps gives the
claim. Resource declarations therefore concern actual shared objects, including
random-number streams and external callbacks; type compatibility alone does
not establish this property.

\paragraph{Counterexample and positive control.}
Starting at $s=1$, two independent data branches that perform $s\leftarrow s+1$
and $s\leftarrow 2s$ yield $4$ or $3$ under opposite topological orders.
Declaring both branches as writers of the same resource introduces the chosen
control order and preserves the corresponding result in serial and threaded
execution. The test suite also uses synchronization barriers to verify that
independent branches actually overlap in time. It exercises read/write
conflicts, wildcard resources, zero-output actions, graph edits, storage
reordering, and failure propagation.

\paragraph{Training and persistence.}
Candidates are prepared before evaluation. An accepted candidate retains the
evaluated graph and its semantic order; method replacement preserves stable
node identities. Graph growth, replacement, and subprogram extraction validate
the resulting dependencies. Explicitly changing the order creates a new
candidate to evaluate. Saved plans contain the order, ports, method signatures,
and effect declarations; loading checks these contracts before replacing the
live model. The plan stores execution semantics rather than a copy of the
external environment.

\paragraph{Optional execution mode.}
V1 enables effect ordering by default and records the choice with the model.
With ordering enabled, a stateful batch completes one sample's graph before
the next sample begins. With ordering disabled, scheduling uses data edges
alone and retains node-major batch processing. Both modes validate graph ports
and preserve the chosen setting through training, cloning, and persistence.
Switching modes changes the program's execution contract and calls for a new
evaluation. The ordering guarantee applies to successful executions; external
effects already performed before a method failure are not rolled back.

\paragraph{Frozen-model regression.}
All 90 final models from the type experiment were replayed in V1 on their 20
initial conditions. Across 1,800 model--initial-state pairs and 4,934 paired
environment steps, serial and three-worker execution matched in trajectories,
accuracy, and loss, including the saved V0 outcomes. The associated implementation
suite passes 85 tests and 46 subtests. These results validate the integrated
execution path without repeating model training.

\section{Effect of Type Resolution on Stateful Search}
\label{app:revision_types}

Separating the roles of position and target produces the strongest search result
in the finite attractor task: 29 of 30 runs solve every initial condition.
Merging those roles lowers success to 15 of 30, and using scalar types throughout
lowers it to 1 of 30. The experiment directly exposes the contribution of the
supplied type prior.

\paragraph{Protocol.}
The three conditions share the nine original method implementations, arities,
groups, weights, and registration order. Fine types distinguish
\texttt{pos\_type} and \texttt{target\_type}. Coarse types merge both into
\texttt{quantity}, which remains distinct from \texttt{scalar}. Flat types use
\texttt{scalar} throughout. Every condition retains the goal-aware movement
operator. Formal seeds are 110600--110629; pilot seed 110500 is excluded.
Python and NumPy random generators are reset to the run seed for each condition.

GraphEvo uses three fixed-mode initial models, three initialization layers,
30 refinement steps per cycle, five maximum evolution cycles, the full
candidate pool, and no method-pool evolution. Early stopping is enabled at zero
loss. Every evaluation exhaustively visits the 20 ordered pairs
$(p,g)\in\{-2,-1,0,1,2\}^2$ with $p\ne g$, resetting the external position and
target for each episode. A graph receives at most four executions to reach the
target. Loss is mean absolute final distance; a successful run reaches the
target from all 20 initial states. Evaluation performs no random sampling and
checks that it leaves the search random-number states unchanged.

\begin{table}[tb]
\centering
\small
\caption{Type resolution under the frozen V0 search protocol. Wilson intervals
describe the run-level success proportion. Evaluation counts are means per run.}
\label{tab:revision_types_full}
\begin{tabular}{lrrrr}
\toprule
Type condition & Successes & 95\% interval & Evaluations & Distinct graphs \\
\midrule
Fine & 29/30 & 83.3--99.4\% & 59.47 & 6.87 \\
Coarse & 15/30 & 33.2--66.8\% & 81.67 & 5.60 \\
Flat & 1/30 & 0.6--16.7\% & 738.67 & 22.77 \\
\bottomrule
\end{tabular}
\end{table}

\paragraph{Search progression and accounting.}
Initialization alone succeeds in 20, 4, and 0 runs, respectively. Subsequent
refinement raises these totals to 29, 15, and 1. Mean final accuracies are
0.9667, 0.5033, and 0.3200, with mean losses 0.0667, 0.9900, and 1.0667.
Nominal caps are shared; early stopping and the available replacements produce
the different actual evaluation budgets in Table~\ref{tab:revision_types_full}.
``Distinct graph'' means an ordered attributed-graph signature, with literal
display names removed; it is an accounting unit rather than a functional
equivalence class.

The V0 early-stop path exits before updating some refinement counters and the
best-model snapshot. We therefore evaluate the saved final graph and record
every evaluation callback. Reconstructing accepted replacements from these
callbacks gives totals of 9, 13, and 19, compared with native logged counters
of 0, 1, and 18. This bookkeeping correction changes neither the trained graphs
nor their measured outcomes. The earlier ten-seed experiment in
Appendix~\ref{sec:appendix_exp4} uses a different episode-sampling protocol and
remains a separate experiment. The present result establishes improved search
success from a role-sensitive type prior in this fully enumerated task.

\section{Matched Explicit and Implicit State Encodings}
\label{app:revision_encoding}

Explicit state threading and implicit external-state access reproduce the same
program behavior when they preserve the same state-access order. We verified
this correspondence across all exported records from the type experiment,
including unsuccessful candidates.

\paragraph{Construction.}
For each run and condition, the candidate catalogue contains the first 20
distinct ordered graph signatures encountered, where available, plus the final
graph even when already recorded. The resulting 1,007 records include repeated
final graphs. Every state-reading or state-writing method is translated to
consume and return an immutable $(p,g)$ state token. Token edges follow the
original serial state-access order, including reads. Each access adds one state
input/output port and one state edge; a final edge returns the resulting state.
No computation node is added. Numerical argument binding and output aggregation
remain unchanged.

\paragraph{Execution comparison.}
Both encodings use the original native serial executor on the same 20 initial
conditions and four-step episode limit. Checks compare per-node numerical
outputs, state after each access, collapsed output, each environment transition,
and the recorded final loss and accuracy. During explicit execution, the global
implicit state is deliberately assigned unrelated values, which verifies that
the token-based methods neither read nor modify it.

\begin{table}[tb]
\centering
\small
\caption{Matched state encoding results using V0. Ratios include the common
adapter interface and state-token construction/extraction.}
\label{tab:revision_encoding_full}
\begin{tabular}{lr}
\toprule
Quantity & Result \\
\midrule
Candidate/final-model records & 1,007 \\
Paired initial conditions & 20,140 \\
Paired environment steps & 70,670 \\
Native graph executions, both encodings & 141,340 \\
Numerical, state, trajectory, and metric comparisons & All matched \\
Mean additional state edges, final graphs & 2.189 \\
Median explicit/implicit elapsed-time ratio & 1.137 \\
Ratio interquartile range & 1.086--1.197 \\
\bottomrule
\end{tabular}
\end{table}

\paragraph{Timing.}
Each of the 90 final graphs receives one complete warm-up and five paired timing
repetitions over all 20 episodes; the order of the two encodings is randomized
within each pair. The median ratio is computed per model and then summarized
across models. It measures this Python object-port implementation, including
state construction, dispatch, and extraction. Both conditions use the same
already-searched candidates, so the experiment isolates encoding and execution;
independent search over explicit-state programs is a separate comparison.

\section{Controlled Single-Graph Execution Scaling}
\label{app:revision_scaling}

At fixed depth and batch size, the native executor processes graphs from 32 to
1,024 operations with consistent numerical results. Over this 32-fold increase
in operation count, median serial time increases from 2.127 to 70.586\,ms.
Type compatibility and topological preparation remain small relative to graph
execution throughout the measured range.

\paragraph{Graph and workload design.}
Graphs have eight layers of width $w\in\{4,8,16,32,64,128\}$, four input
features, and batch size 32. There are $8w$ operation nodes, two boundary nodes,
and $2(8w)+w$ edges. Each operation has two inputs. Its first predecessor slot
uses a permutation of the preceding layer, ensuring every earlier node
contributes to an output, while the second slot selects a distinct predecessor.
The first layer selects two of the four input features. Final-layer outputs
are aggregated with the native average operation.

Each scalar operation starts from $z=0.6a+0.4b$ and performs 16 iterations of
$z\leftarrow0.98z+0.02\tanh(0.9z+0.05)$. This is ordinary CPU arithmetic,
with vectorization disabled. Five independently wired graphs are evaluated at
each width. Two disjoint pilot graphs precede the frozen formal protocol.
The input seed is 2026110600, and formal graph seeds at width $w$ are
$51000+1000w+i$, $i\in\{1,\ldots,5\}$.

\begin{table}[tb]
\centering
\small
\caption{V0 execution scaling, in milliseconds for batch size 32. Values are
medians over five graph-level timing medians. Speedup is the median of paired
serial/threaded ratios, so it need not equal the ratio of the displayed medians.}
\label{tab:revision_scaling_full}
\begin{tabular}{rrrrrrr}
\toprule
Operations & Edges & Type pass & Topology & Serial & 4 threads & Speedup \\
\midrule
32 & 68 & 0.051 & 0.090 & 2.127 & 3.478 & 0.610 \\
64 & 136 & 0.096 & 0.164 & 4.086 & 5.782 & 0.697 \\
128 & 272 & 0.181 & 0.305 & 8.022 & 10.482 & 0.754 \\
256 & 544 & 0.367 & 0.598 & 16.894 & 20.436 & 0.805 \\
512 & 1,088 & 0.753 & 1.163 & 32.621 & 40.360 & 0.808 \\
1,024 & 2,176 & 1.448 & 2.310 & 70.586 & 82.379 & 0.811 \\
\bottomrule
\end{tabular}
\end{table}

\paragraph{Measurement and validation.}
Each metric receives two warm-ups and seven timed calls, with shuffled case and
metric order. Timing uses a monotonic high-resolution clock without memory
tracing. The environment is Windows, Python 3.12.9, NumPy 2.1.3, and NetworkX
3.4.2, on a 16-core/24-thread Intel host. Serial timing includes topological
sorting; threaded timing includes dependency bookkeeping and thread-pool
creation/destruction. The separate topology column is not deducted from either.
The type pass is an additional per-edge measurement calling the native type
compatibility predicate on declared ports, including the final aggregation
contract. It is distinct from an integrated whole-program compiler check.

All 30 graphs satisfy the specified counts, depth, acyclicity, input slots, and
declared type compatibility. A deliberately incompatible signature is rejected.
Both execution paths agree with an independently implemented NumPy reference
at tolerance $10^{-12}$; maximum observed reference error is
$2.22\times10^{-16}$. Separately traced allocation peaks exclude pre-existing
graphs and inputs, and measure Python allocations rather than total resident
memory. At 1,024 operations, the median serial and threaded peaks are 3.60 and
5.33\,MB, respectively.

\paragraph{Observed scaling regime.}
Four Python threads are slower than serial execution for this scalar CPU
workload. Together with the earlier simulated I/O experiment, the measurements
show how operator cost and GIL behavior determine whether available graph
parallelism improves latency. The largest tested graph has 1,024 operations;
the point at which preparation dominates has not been reached. Node and edge
counts grow together at fixed arity, so this experiment measures a controlled
width expansion rather than an independent edge-density sweep.

\section{Shared-Feature Graphs for Action Selection}
\label{app:revision_action}

A shared-feature graph provides a concrete connection between numerical
computation, persistent state, and discrete actions. One feature can contribute
to several action scores while a mutually exclusive routing stage selects a
zero-output action node. The same organization supports the four evaluated
task families.

\paragraph{Controller representation and optimization.}
For observation $o_t$ and previous internal state $h_t$, the graph computes
\begin{equation}
 z_t=\phi_G(o_t,h_t),\qquad
 q_{t,a}=\sum_j w_{a,j}z_{t,j},\qquad
 a_t=\operatorname{route}(q_t,u_t),
\label{eq:revision_action_model}
\end{equation}
where $u_t$ supplies any stochastic routing draw. Feature calculations read the
old state, and a single commit stage writes the state used on the next step.
The shared configuration has four storage slots arranged as two independent
state channels. A decay factor of 0.9 and clipping are part of its specified
update rule.

The training procedure provides the shared-feature interface, per-action score
branches, route selection, and state-update rule. CMA optimizes the score
coefficients. Configurations permitting growth propose eight additional
arithmetic features at each of the prescribed 15\%, 35\%, 55\%, and 75\%
budget milestones and retarget state-emission sources. Thus learned coefficients
and proposed features sit within a declared controller template. Models are
selected by validation success, with progress breaking ties. Each task uses
its own selected model; the task identity is supplied externally.

\begin{table}[tb]
\centering
\small
\caption{Held-out completion rates (\%) of the selected controller portfolio.
Each cell uses 2,048 episodes; all five training seeds are shown. The mean row
averages training seeds within each task.}
\label{tab:revision_action_seeds}
\begin{tabular}{rrrrr}
\toprule
Training seed & Catch & Maze & Snake & Bird \\
\midrule
101 & 99.22 & 96.09 & 91.65 & 100.00 \\
102 & 61.28 & 75.98 & 0.78 & 100.00 \\
103 & 100.00 & 98.88 & 32.18 & 100.00 \\
104 & 100.00 & 96.24 & 93.12 & 100.00 \\
105 & 100.00 & 96.53 & 0.05 & 100.00 \\
\midrule
Mean & 92.10 & 92.74 & 43.55 & 100.00 \\
\bottomrule
\end{tabular}
\end{table}

The bird results are consistent across seeds, while Snake shows substantial
variation. In the validation-selected representative system (seed 104), Catch,
Maze, and Snake have no subsequent feature-growth event; Bird has four. Its
controller contains 44 features and 88 coefficients. One saved feature is
$z_{16}=\mathrm{gapBottom}+h_{1,t}$ and contributes with coefficients about
$+3.001$ and $-6.538$ to the flap and coast scores. This is a direct example of
shared numerical computation feeding distinct action branches.

\paragraph{Execution on the original ATF backend.}
The frozen representative Bird graph has 314 nodes and 411 edges. An explicit
adapter expands the input, constant-trigger, state-commit, and observation
interfaces, giving 316 nodes and 565 edges. Every feature operation, score
coefficient, product, and accumulation is represented by graph operations.
The commit stage sends a completion token to the router; enabled zero-output
action sinks append action receipts to the step context, and the environment
advances after graph execution.

We replay the first 20 consecutive episodes of the pre-existing test stream
(environment seed 200000104; episode indices starting at 83000000), using the
same external random draw in all executors. Across 20,480 steps, V0 serial
execution agrees exactly in features, scores, state, actions, and trajectories
with the portable interpreter and the original kernel function body. A further
2,048 of these steps match the previously archived trace. All 49 boundary
probes and their 49 insertion-order-reordered counterparts pass. This replay
establishes executable transfer of the frozen controller through the stated
adapter; the performance results in Table~\ref{tab:revision_action_seeds}
come from the preceding training and held-out evaluation.

\section{Sensitivity to the Atomic Method Pool}
\label{app:revision_pool}

Expanding the atomic pool changes search quality non-monotonically. The best
aggregate numerical result in this comparison comes from the six-operation
core, while the best maze result comes from the nine-operation extension.
The comparison covers nine pools spanning 6--20 operations.

\paragraph{Pools and search protocol.}
The core contains identity, addition, subtraction, multiplication, safe
division, and duplication. The nested extensions add increment, decrement,
and negation (9 operations); square, ReLU, absolute value, three-input mean,
maximum, and minimum (15); then sine, cosine, two-argument arctangent, scaling,
and safe square root (20). Five further pools vary polynomial, reciprocal,
and trigonometric combinations. All use the same path-feedback procedure:
each distinct executed route receives one terminal-loss update per episode,
with exponential averaging coefficient 0.08, temperature 0.12, and exploration
fraction 0.2. Candidates grow from identity scaffolds with periodic extension.

Numerical tasks use 131,072 training input presentations per task and seed,
sampled from 2,048 fixed training examples, 256 candidate opportunities, and a
32-node limit per graph. Twelve tasks are grouped into ten equally weighted
families, averaging the three equivalent two-variable multiplication tasks
before aggregation. Four training seeds are used. Testing uses 1,024 inputs;
the reported pass rate is the fraction with relative error at most 1\%.
Action tasks use 65,536 training episodes per task and seed, 128 candidate
opportunities, a 24-node graph limit, and two seeds. Checkpoint evaluation uses
512 episodes per route-execution mode. Candidate generation, pool composition,
and path feedback define this subsequent search host.

\begin{table}[tb]
\centering
\small
\caption{Atomic-pool sensitivity. Numerical values average ten task families
and four seeds. Maze values average two seeds under within-episode route
switching. Both columns are percentages with their own evaluation units.}
\label{tab:revision_pool_full}
\begin{tabular}{lrrr}
\toprule
Pool & Operations & Numerical input pass rate & Maze completion \\
\midrule
Arithmetic & 6 & 14.01 & 11.43 \\
Square & 7 & 11.47 & 17.97 \\
Basic & 9 & 11.45 & 58.01 \\
Nonlinear & 15 & 7.72 & 37.89 \\
Extended & 20 & 11.35 & 6.93 \\
Algebra & 8 & 11.32 & 12.89 \\
Polynomial & 9 & 8.57 & 31.84 \\
Balanced & 12 & 13.62 & 27.44 \\
General & 16 & 10.83 & 14.16 \\
\bottomrule
\end{tabular}
\end{table}

The numerical comparison comprises 432 records, including 32 reused baseline
records; the action comparison comprises 144 newly trained runs. The broader
928-run study additionally includes mechanism ablations, so that total is not
the number of pool-size comparisons. More operations offer new expressions
while increasing the alternatives explored within the same nominal budget.
Because operator identities and arities change with the extensions, the table
describes practical pool sensitivity. Numerical input pass rates and action
completion rates are reported separately, preserving their distinct meanings.

\section{Large Continual-Learning Execution Workload}
\label{app:revision_large}

The ATF execution core also supports a sustained workload with 80 complete
learning sequences and 3,200 stages. Across training and checkpoint evaluation,
239,206,400 input rows are checked against native ATF outputs. This measures
the substrate's use inside a long-running application, complementing the
single-graph measurements in Appendix~\ref{app:revision_scaling}.

\paragraph{Environment and design.}
Each observation has 64 fields, one of 32 identities, and six available actions.
The 40-stage sequence first presents 20 deterministic rule families and then
returns to them in reverse order. Fields and identities are assigned by seed.
Training uses 3,072 interactions per stage, at most 2,048 candidates per search
refresh, 6,144 stored observations, and a 48-derived-node cap per candidate
graph. Ten seeds are used:
41113, 41213, 41317, 41423, 41527, 41631, 41741, 41847, 41957, and 42061.

Six mechanisms are compared at concept-library capacity 32: direct reuse,
small-set retrieval, retrieval with prediction history, reserved fresh search,
reward/complexity selection, and their full combination. Direct reuse and the
full combination are additionally run at capacity 128, yielding 80 sequences.
Paired conditions share observations, exploration draws, action sets, and
nominal training and candidate limits. ATF executes the risk-expression graph,
identity lookup, and action-value computation. Candidate proposals, selection,
and library management are supplied by the surrounding learner.

\paragraph{Counting execution.}
The completed runs contain 9,830,400 effective training actions and 229,376,000
checkpoint-evaluation rows. Their sum is the 239,206,400 native-checked input
rows reported above. One graph call can receive a NumPy array containing many
rows, and the action-value head can perform a matrix computation within one
node. Consequently, input rows, graph calls, concept entries, and graph nodes
are distinct scale measures. The independent training units remain the ten
seeds.

\paragraph{Learning outcomes.}
Learning cost is the earliest checkpoint after which both balanced risk accuracy
of at least 90\% and optimal-action rate of at least 80\% in the carrying
condition are maintained. Unsuccessful stages receive the fixed cost 3,200.
A rule/first-or-return configuration counts as mastered when at least eight
of ten seeds meet the criterion, giving 40 configurations in total.

\begin{table}[tb]
\centering
\small
\caption{Continual-learning results across ten paired seeds. Lower cost and
training loss are better. Loss is accumulated reward shortfall relative to an
oracle action over the 40 stages.}
\label{tab:revision_large_results}
\begin{tabular}{lrrrr}
\toprule
Mechanism (capacity) & Mean cost & Return cost & Training loss & Mastered/40 \\
\midrule
Direct reuse (32) & 1,792.32 & 1,749.76 & 62,809.12 & 21 \\
Small-set retrieval (32) & 1,777.60 & 1,762.56 & 62,942.44 & 21 \\
Retrieval + history (32) & 1,720.32 & 1,648.00 & 62,617.83 & 23 \\
Reserved fresh search (32) & 1,817.92 & 1,739.52 & 63,083.45 & 21 \\
Reward/complexity (32) & 1,800.64 & 1,767.68 & 62,681.52 & 21 \\
Full combination (32) & 1,855.04 & 1,802.88 & 63,203.22 & 20 \\
Direct reuse (128) & 1,691.52 & 1,512.96 & 61,926.77 & 21 \\
Full combination (128) & 1,683.20 & 1,472.00 & 61,921.29 & 21 \\
\bottomrule
\end{tabular}
\end{table}

The prespecified advantage criterion for the full combination at capacity 32
requires at least eight of ten seeds to reduce cost by at least 10\%, without
increasing average training loss. It is not met: zero seeds attain the cost
threshold, mean cost increases by 3.50\%, and loss increases by approximately
0.63\%. Retrieval with history has the lowest mean cost among the capacity-32
conditions. Higher-order parity remains difficult: the three- and five-input
XOR rules have approximately chance-level risk accuracy, with no seed reaching
the mastery criterion on first or return exposure.

\paragraph{Completion audit.}
An initial 16\,MiB serialized-state guard stopped part of the experiment.
The formal completion uniformly raises that guard to 256\,MiB while retaining
the training, candidate, and concept budgets. Fifty-two completed sequences
are preserved with identical file hashes; the remaining 28 start from their
initial conditions. Three discarded prefixes account for 343,296 additional
actions, recorded separately from the effective training count. Final frozen
replay passes for all 80 models. Peak serialized state for the full combination
is approximately 4.67\,MiB at capacity 32 and 18.67\,MiB at capacity 128.

\section{Reusable-Library Capacity and Return to Earlier Tasks}
\label{app:revision_capacity}

Larger libraries help most when earlier tasks must be recovered after a longer
sequence. In the 40-stage experiment, increasing capacity from 32 to 128 lowers
mean cost by 5.62\% for direct reuse and 9.26\% for the full combination.
For the full combination, return-stage cost drops from 1,802.88 to 1,472.00
($18.35\%$), whereas first-exposure cost changes from 1,907.20 to 1,894.40.
At the same capacity 128, the full combination and direct reuse are close:
their mean cost difference is $-8.32$, with an exploratory paired-bootstrap
95\% interval of $[-103.68,87.68]$. The larger-storage result therefore concerns
capacity and reuse over time.

\paragraph{Paired short-sequence comparison.}
A complementary experiment changes only the library capacity to 8, 16, 32, or
64, retaining the learner, tasks, admission and replacement rules, search
budgets, and five training seeds. Three five-stage relationship sequences per
seed give 60 sequences, 300 stages, and 460,800 interactions. Transfer cost
averages the new-identity and return stages; a failure receives the fixed cost
1,600. This is a paired capacity sensitivity test on the same seeds.

\begin{table}[tb]
\centering
\small
\caption{Short-sequence library capacity comparison. The actual peak counts
show whether the configured limit was reached. All capacities master 14 of the
15 sequence-stage configurations under the frozen criterion.}
\label{tab:revision_capacity_short}
\begin{tabular}{rrrrr}
\toprule
Capacity & Transfer cost & Training loss & Mastered/15 & Peak entries \\
\midrule
8 & 226.13 & 2,993.39 & 14 & 8 \\
16 & 238.93 & 3,000.39 & 14 & 16 \\
32 & 234.67 & 2,997.98 & 14 & 21 \\
64 & 234.67 & 2,997.98 & 14 & 21 \\
\bottomrule
\end{tabular}
\end{table}

Increasing capacity in these short sequences produces no stable learning gain.
All three larger-capacity conditions have zero of five seeds achieving the
prespecified 20\% transfer-cost improvement. Capacities 32 and 64 generate
identical models, learning curves, libraries, and action trajectories in all
15 paired sequences because occupancy never exceeds 21 entries. Meanwhile,
the mean final single-model derived-node count is 15.87, 15.20, 15.19, and 15.19.
Library capacity controls retained reusable structures; the final selected
graph can remain small even when the collection of stored structures grows.

\paragraph{Evidence organization.}
The four V0 supplement records are organized under the relative artifact
directories \path{experiments/rebuttal_20261106/type_search},
\path{state_encoding}, \path{scaling}, and \path{action_example}; each contains
its protocol, source-version record, and measurements. The subsequent studies
are identified by their controller portfolio, atomic-pool sensitivity,
40-stage continual-learning, and short-sequence capacity protocols. Their
source snapshots and row-count audits distinguish training, evaluation, and
diagnostic records. The results above use only completed formal records;
planned extensions and withdrawn records do not enter the tables.
